
\SGAProof{Proof of \SGARefLink{sga3:xvii:lemma:3:9:item:iii}{3.9(iii)}}
The group \(G\) is étale.  Proceeding as in (ii), we reduce first to
the case in which \(G\) is commutative, and then to the case in which
\(G\) is annihilated by \(p\).  But then
\[
  G_{\overline{k}}\simeq(\mathbb Z/p\mathbb Z)^r.
\]

\SGAProof{Proof of \SGARefLink{sga3:xvii:theorem:3-5}{3.5}, \SGARefLink{sga3:xvii:theorem:3:5:item:i}{(i)}\(\Rightarrow\)\SGARefLink{sga3:xvii:theorem:3:5:item:iv}{(iv)}, in case \SGARefLink{sga3:xvii:section:b}{(b)},
\(p=0\)}
The group \(G\) is then smooth and connected; proceeding as in
\SGARefLink{sga3:xvii:lemma:3:9:item:ii}{3.9(ii)}, we reduce to the case in which \(G\) is also commutative.
We then have the following more precise result.

\SGAStatement{Lemma 3.9 ter}{Let \(k\) be a field of characteristic
\(0\), let \(G\) be a commutative unipotent \(k\)-algebraic group,
and put \(\mathfrak g=\operatorname{Lie}G\).  There is a canonical
isomorphism
\[
  \exp:\mathbf W(\mathfrak g)\xrightarrow{\ \sim\ }G.
\]
The morphism \(\exp\) is the unique homomorphism
\(\mathbf W(\mathfrak g)\to G\) that induces the identity on Lie
algebras.}
\SGARefTarget{sga3:xvii:lemma:3:9:ter}

Since \(G\) is unipotent, it can be realized as an algebraic
subgroup of \(\operatorname{Trigstr}(n)\), for a suitable integer
\(n\) (\SGARefLink{sga3:xvii:theorem:3-5}{3.5}, \SGARefLink{sga3:xvii:theorem:3:5:item:i}{(i)}\(\Rightarrow\)\SGARefLink{sga3:xvii:theorem:3:5:item:v}{(v)}).  Choosing such an embedding
identifies \(\mathfrak g\) with a Lie subalgebra of
\(\mathfrak{gl}(n)\) consisting of nilpotent endomorphisms.  We
thereby obtain a \(k\)-morphism
\[
  \begin{aligned}
  \exp:\mathbf W(\mathfrak g)&\longrightarrow
    \operatorname{GL}(n),\\
  T&\longmapsto\sum_{i\geq0}\frac{T^i}{i!}.
  \end{aligned}
\]

Since \(G\) is commutative, the morphism \(\exp\) is a
homomorphism.  Let \(G'\) be the algebraic group that is the image
of \(\mathbf W(\mathfrak g)\) under \(\exp\).  Under the canonical
identification
\(\operatorname{Lie}\mathbf W(\mathfrak g)=\mathfrak g\), the
tangent linear map of \(\exp\) is simply the inclusion
\(\mathfrak g\hookrightarrow\mathfrak{gl}(n)\).  Hence
\[
  \operatorname{Lie}(G\cap G')
  =\mathfrak g\cap\mathfrak g=\mathfrak g.
\]
Since \(G\cap G'\) is smooth (Exposé VI\textsubscript{B}, \SGARefLink{cor:sga3-VIB-1.6.1}{1.6.1})
and \(G\) is connected---being an iterated extension of groups
\(\mathbb G_a\) (\SGARefLink{sga3:xvii:theorem:3-5}{3.5}, \SGARefLink{sga3:xvii:theorem:3:5:item:i}{(i)}\(\Rightarrow\)\SGARefLink{sga3:xvii:theorem:3:5:item:ii}{(ii)})---one necessarily has
\(G=G'\).  The kernel \(\ker(\exp)\) is an étale unipotent group,
and hence is the trivial group.  Consequently \(\exp\) is an
isomorphism from \(\mathbf W(\mathfrak g)\) onto \(G\).

If \(h:\mathbf W(\mathfrak g)\to G\) is another homomorphism such
that \(\operatorname{Lie}(h)\) is the identity map of
\(\mathfrak g\), then \(h-\exp\) is a homomorphism, since \(G\) is
commutative, and its tangent linear map is zero.  Because \(k\) has
characteristic \(0\) and \(\mathbf W(\mathfrak g)\) is connected,
Cartier's theorem again implies that \(h=\exp\).
